% Chapter 1

\chapter{\uppercase{Introduction}}
\label{ch:intro}
AI agents now work on coding projects that last weeks and span many sessions. Each new session
starts with an empty prompt, so the user has to repeat what they said before. This project
gives the agent a long-term memory that lives outside the model.

\section{\uppercase{Background}}
Pasting the whole history into every prompt does not work well. Models miss details in the
middle of long inputs \cite{liu2024lost}, and old statements sit next to their corrections with
no rule for which is current. So we built a separate memory service for the Pi coding agent
\cite{zechner2026pi}. It stores the original messages, extracts assertions from them and records
when an assertion changed. For a question it returns only the relevant assertions and their
quotes. Here \emph{memory} means this stored knowledge, not RAM; other terms are in the Glossary at the end of the report.

The engine has five modules. M1 Source Acquisition and Processing saves the messages. M2
Knowledge Extraction and Consolidation extracts assertions. M3 Temporal Knowledge and Graph
Management stores them with their history. M4 Hybrid Retrieval and Grounded Reasoning finds
the right ones for a question. M5 Model Training and Evaluation trains the small local model
the others use. M1 to M4 form the memory service; M5 runs offline.

\section{\uppercase{Problem Statement}}
Agents without persistent memory lose instructions and decisions between sessions. When an
instruction changes, they may reuse the old value, so users repeat context and fix the same
mistakes.

\section{\uppercase{Objectives}}
\begin{enumerate}[itemsep=0pt]
\item Retain useful information for AI agents across sessions.
\item Keep the updated instruction while preserving earlier versions.
\item Link every stored assertion to the original message so it can be verified.
\item Fine-tune a small local model to extract assertions, combine them and answer with citations.
\item Retrieve relevant information quickly within a limited context size.
\item Evaluate the engine on a public long-conversation benchmark, LoCoMo \cite{maharana2024locomo}.
\end{enumerate}

\section{\uppercase{Challenges Faced}}
Our early models did well on generated data but passed only 18 of 33 real conversation
segments. They mixed up quotes, sources and dates, so we revised the data, prompts and checks
several times. The hardest bug was that the deployed model silently cut off long inputs. We now
check the input length and split long conversations.

\section{\uppercase{Organisation of the Report}}
Chapter~\ref{ch:survey} reviews existing memory systems. Chapter~\ref{ch:design} describes the
architecture and the modules with one running example. Chapter~\ref{ch:impl} explains how each
module is built. Chapter~\ref{ch:results} covers training, a demo and the benchmark results.
Chapter~\ref{ch:conclusion} gives the conclusion, limitations and future work.
Appendix~\ref{app:schema} lists the database schema and Appendix~\ref{app:prompts} the prompts of the local model.
